\chapter{Fazit}
Der erste Teil der Zielsetzung dieser Arbeit bestand darin, die mathematischen Eigenschaften elliptischer Kurven zu untersuchen und ihre Bedeutung für die Public-Key-Krpytographie darzustellen. Hierbei wurde gezeigt, dass die Punkte einer elliptischen Kurve zusammen mit der Punktaddition eine abelsche Gruppe bilden. Mithilfe der Additionsformeln lassen sich Gruppenoperationen durchführen, wobei insbesondere die Skalarmultiplikation durch Verfahren wie Double-and-Add effizient berechnet werden kann.

Anschließend wurden elliptische Kurven über endlichen Körper $\mathbb{F}_q$ betrachtet. Die endliche Gruppenstruktur ermöglicht ebenfalls eine abelsche Gruppe. Dadurch kann effizienter und exakt auf einem Computer gerechnet werden, da keine Gleitpunktarithmetik benötigt wird. Ein Element eines endlcihen Körper $\mathbb{F}_q$ kann durch höchstens \(\lceil \log_2 q\rceil\) Bits dargestellt werden. Dadurch bleibt die Größe der Darstellung der Körperelemente beschränkt. Die endliche Gruppenstruktur ermöglicht es, dass diskrete Logarithmus Problem auch auf elliptischen Kruven zu formulieren. Hierbei ist ein Punkt $P$ und ein Punkt $Q=kP$ gegeben, dass Skalar $k$ ist hierbei unbekannt. Für kryptographisch geeignete Kurven sind nach aktuellem Kenntnisstand keine klassischen Algorithmen bekannt, die das allgemeine elliptische diskrete Logarithmusproblem in Polynomialzeit lösen.

Auf dieser Eigenschaft des elliptischen diskreten Logarithmusproblems lassen sich kryptographische Verfahren aufbauen. Während die Berechnung von $Q=kP$ mithilfe dem Double-And-Add-Algorithmus effizient möglich ist, gilt die umgekehrte Bestimmung des Skalars \(k\) aus den gegebenen Punkten \(P\) und \(Q\) bei geeigneter Parameterwahl nach aktuellem Kenntnisstand als rechnerisch schwierig.

Ein zentrales mathematisches Ergebnis stellt der Satz von Hasse dar
\[
|\#E(\mathbb{F}_q)-q-1 |\leq 2\sqrt q
\]
welcher die Anzahl der rationalen Punkte einer elliptischen Kurve über \(\mathbb F_q\) einschränkt. Dieser leifert eine wichtige Grundlage für die Untersuchung der Gruppenordnung. Der Schoof-Algorithmus ermöglicht es, die Gruppenordnung über endlichen Körpern effizient und exakt zu bestimmen.

Neben den mathematischen Grundlagen wurden auch verschiedene klassische Angriffe auf das elliptische diskrete Logarithmusproblem untersucht. Der Baby-Step-Giant-Step-Algorithmus zeigt mit einer Laufzeit von \(\mathcal O(\sqrt n)\) Gruppenoperationen, wobei \(n\) die Ordnung des Basispunktes bezeichnet, dass ausreichend große Untergruppen erforderlich sind. Der Pohlig-Hellman-Algorithmus verdeutlicht die Bedeutung der Primfaktorzerlegung der Gruppenordnung. Das diskrete Logarithmusproblem wesentlich effizienter gelöst werden, wenn die Gruppenordnung ausschließlich aus kleinen Primfaktoren besteht.

Beim MOV-Algorithmus handelt es sich um eine spezielle Methode, die die Eigenschaften von gewissen elliptischen Kurven ausnutzt. Unter gewissen Vorraussetzung lässt sich mithilfe der Weil-Paarung das ECDLP in  $E(\mathbb{F}_q)$ auf das DL-Problem in der multiplikativen Gruppe $\mathbb{F}_{q^m}^\times$ übertragen. Dabei spielt insbesondere der Einbettungsgrad eine entscheidende Rolle. Supersinguläre Kurven besitzen in den betrachteten klassischen Anwendungsfällen häufig einen kleinen Einbettungsgrad und können deshalb gegenüber entsprechenden Angriffen besonders anfällig sein.

Ein Vorteil von ellitpischen Kurven ist, dass eine deutlich kürzere SChlüssellänge als bei RSA benötigt wird und das es schwieriger zu lösen ist als das DL-Problem auf der multiplikativen Gruppe, da für diese subexponentielle Angriffe exestieren, welche für elliptische Kurven keine entsprechenden Verfahren bekannt sind.

Es wurde hierbei nur eine Auswahl an Angriffen und ihren Konsequenzen betrachtet. Für praktische Anwendungen empfiehlt sich deshalb die Verwendung etablierter, öffentlich untersuchter und aktuell empfohlener Kurvenparameter.  Damit konnte auch der zweite Teil der Zielsetzung erfüllt werden.

Abschließend wurden die Auswirkungen von Quantencomputern auf die Sicherheit kryptographischer Verfahren auf elliptischen Kurven untersucht. Shors Algorithmus lässt sich auf das elliptische diskrete Logarithmusproblem übertragen und ermöglicht dessen Lösung in Polynomialzeit bezüglich der Bitlänge der verwendeten Untergruppenordnung. Ein hinreichend leistungsfähiger Quantencomputer könnte damit die kryptographische Verfahren auf Grundlage elliptischer Kurven brechen. Dies betrifft auch Kurven mit großer Primzahlordnung der verwendeten Untergruppe und großem Einbettungsgrad. Die zuvor aus klassischen Angriffen abgeleiteten Anforderungen an die Kurvenparameter bieten daher keinen ausreichenden Schutz gegenüber Shors Algorithmus.


Größere Parameter erhöhen zwar auch den Ressourcenbedarf eines Quantenangriffs und können dessen praktische Durchführung erschweren. Während der Aufwand klassischer generischer Angriffe jedoch exponentiell mit der Bitlänge der Untergruppenordnung wächst, bleibt der Aufwand von Shors Algorithmus polynomial. Eine bloße Vergrößerung der Parameter stellt deshalb keine grundlegende Lösung für diese Bedrohung dar. Da der Zeitpunkt der praktischen Realisierung kryptographisch relevanter Quantencomputer weiterhin ungewiss ist, ihre zukünftige Verfügbarkeit jedoch ein ernst zu nehmendes Sicherheitsrisiko darstellt, gewinnt die frühzeitige Umstellung auf Post-Quanten-Kryptographie zunehmend an Bedeutung.

Anküpfend an diese Arbeite könnte untersucht werden inwiefern elliptische Kurven auch in der Post-Quanten-Kryptographie eine Rolle spielen können. Ein möglicher Ansatz sind isogeniebasierte Verfahren. Isogenien wurden bereits im Zusammenhang mit dem Beweis der Hasse-Schranke als Gruppenhomomorphismen zwischen elliptischen Kurven eingeführt
\[
        \phi:E_1\to E_2
        \qquad\text{mit}\qquad
        \phi(\mathcal O)=\mathcal O.
\]
Die Arbeit von Mishra et al.~\cite{mishra26} gibt einen Überblick über verschiedene isogeniebasierte kryptographische Verfahren. Ein konkretes Beispiel ist der Schlüsselaustausch CSIDH. Dabei gehen beiden Kommunikationspartner von einer gemeinsamen öffentlichen Kurven $E_0$ aus. Jeder Teilnehmer wält einen gehmein Schlüssel, der eine bestimme Gruppenwirkung auf elliptischen Kurven festlegt. Daraus werden die öffentlichen Kurven $E_A$ und $E_B$ berechnet. Aufgrund der Kommutativität der Klassenwirkung können beide Teilnehmer anschließend eine gemeinsame Kurve $E_{AB}$ bestimmen. Aus dieser kann ein gemeinsamer Schlüssel abgeleitet werden. Das Verfahren ähnelt somit dem bereits betrachteten Diffie-Hellman-Schlüsselaustausch. Ein möglicher Gegenstand zukünftiger Untersuchungen wären die mathematischen Grundlagen solcher Verfahren sowie die Frage, unter welchen Umständen ein konkretes isogeniebasiertes Verfahren gegenüber klassischen und quantenalgorithmischen Angriffen als sicher gelten kann.